\documentclass{article}
\usepackage[utf8]{inputenc}
\usepackage[T1]{fontenc}
\usepackage{geometry}
\geometry{
    a4paper,
    left=2.5cm,
    right=2.5cm,
    top=2.5cm,
    bottom=2.5cm
}
\usepackage{graphicx}
\usepackage{float}
\usepackage{hyperref}
\usepackage[english]{babel}
\usepackage{amsmath}
\usepackage{amssymb}
\usepackage{array}
\usepackage{booktabs}
\usepackage{caption}
\usepackage{subcaption}
\usepackage{enumitem}
\usepackage{fancyhdr}
\usepackage{titlesec}
\usepackage{url}
\usepackage{textcomp}
\usepackage{xcolor}
\usepackage{listings}
\usepackage{multirow}
\usepackage{longtable}
\usepackage{tabularx}
\usepackage{times}

% Set up page style
\pagestyle{fancy}
\fancyhf{}
\rfoot{\thepage}
\renewcommand{\headrulewidth}{0pt}

% Title formatting
\titleformat{\chapter}[display]{\normalfont\bfseries}{\chaptertitle\thechapter}{0pt}{\Huge}
\titlespacing*{\chapter}{0pt}{50pt}{40pt}

% Document info
\title{Template Document}
\author{Academic Support Agent}
\date{\today}

\begin{document}
\maketitle

\maketitle


\selectlanguage{spanish}
\begin{abstract}
    \textit{\normalsize El resumen describe de forma concisa y clara el contenido del art culo. Las ideas fundamentales de relevancia, as  como los resultados y las conclusiones generales del trabajo reportado.}
    
    \vspace{0.5cm}

    \textbf{Palabras clave:} Palabras que por s  mismas dan idea del tema o  rea del art culo.
\end{abstract}


\selectlanguage{english}
\begin{abstract}
    \textit{\normalsize The abstract concisely and clearly describes the content of the article. The fundamental ideas of relevance, as well as the results and general conclusions of the reported work.}
    
    \vspace{0.5cm}

    \textbf{Keywords:} Words that by themselves give an idea of the subject or area of the article.
\end{abstract}



\selectlanguage{spanish}
\begin{center}
    \section{Introducci n}
\end{center}

    La introducci n es el texto que sirve para ir induciendo al lector hacia el tema que trata el art culo. Las citas a otras fuentes han de referenciarse con par ntesis cuadrados, como se muestra aqu  \cite{greenwade93}. La numeraci n de dichas referencias se efectuar  de acuerdo al orden de aparici n que estas tengan a lo largo del documento.

    Generalmente, la introducci n finaliza con alguna informaci n relacionada al art culo o con algunas reflexiones sobre el resultado del trabajo que se est  presentando. Al terminar una secci n se dejar n dos renglones en blanco, antes de colocar el t tulo de la siguiente secci n.

    Se pide a los autores que los art culos abarquen un m nimo de 10 p ginas. No se considera un l mite m ximo eso depende del trabajo realizado y lo que se desea difundir.



\newpage % Agrega o elimina estos saltos de p gina de acuerdo a tus necesidades.
\vspace{0.5cm}
\begin{center}
    \section{T tulo}
\end{center}

    El nombre de  ste segundo t tulo depender  de c mo se haya estructurado el art culo.
    
    Se debe considerar una secci n en la cual se presente la metodolog a a seguir, utilizando de preferencia un mapa conceptual (esquemas, diagrama a bloques, de flujo, etc.) para mostrar gr ficamente la metodolog a utilizada. Algunas de las secciones que pudieran llevar la estructura del art culo son: Descripci n del problema, Metodolog a, Modelaci n y Simulaci n, Dise o, Manufactura, Integraci n, Metodolog a, etc.

    \subsection{T tulo de subsecci n}

        Cuando se muestre una ecuaci n en el art culo, dicha ecuaci n ha de enumerarse. La numeraci n se efectuar  de acuerdo al orden de aparici n de la ecuaci n en el art culo. Cuando se haga referencia a una ecuaci n debe hacerse de una manera similar a:  ... como se observa en \ref{eq:pend} , donde \ref{eq:pend} es el n mero de ecuaci n correspondiente. A continuaci n se muestra un ejemplo de una ecuaci n y su numeraci n.

        \vspace{0.5cm}
        \begin{equation}
            m=\frac{y_2-y_1}{x_2-x_1}
            \label{eq:pend}
        \end{equation}
        \vspace{0.5cm}

        Es importante mantener la calidad de la informaci n de los art culos, por esta raz n se sugiere al autor mostrar y describir las figuras que  l considere importantes en su documento. Por ejemplo, en la figura \ref{fig:frog} se muestra una rana.

        De preferencia todas las figuras deben tener su referencia, excepto cuando la figura es producto del trabajo relacionado con el trabajo, o dise ado por los autores.
        
        \vspace{0.5cm}
        \begin{figure}[h!]
        \centering
            \includegraphics[width=0.5\textwidth]{frog}
            \caption{\texttt{\textsl{{\footnotesize Redeyed Tree Frog Smile. [Referencia]}}}}
            \label{fig:frog}
        \end{figure}
        \vspace{0.5cm}

        Se debe agregar un diagrama a bloques que ilustre la metodolog a utilizada en el desarrollo del trabajo, y posteriormente se debe detallar cada uno de los bloques.

        \subsubsection{T tulo de sub-subsecci n}

            En caso de que sea necesario agregar una secci n de tercer nivel, sigue este formato.



\begin{center}
    \section{Resultados}
\end{center}

        Aqu  se muestran los resultados obtenidos, se pueden adjuntar tablas, cit ndolas de la siguiente manera \ref{tab:1}.

        \vspace{0.5cm}
        \begin{table}[h!]
            \centering
                \begin{tabular}{|c|c|c|c|}
                \hline
                \textbf{Caracter stica} & \textbf{1} & \textbf{2} & \textbf{3} \\ \hline
                {Tipo 1}         & x          & x          & x          \\ \hline
                {Tipo 2}         & x          & x          & x          \\ \hline
                {Tipo 3}         & x          & x          & x          \\ \hline
                {Tipo 4}         & x          & x          & x          \\ \hline
            \end{tabular}
            \label{tab:1}
        \end{table}
        \vspace{0.5cm}

        Las tablas se deber n presentar en la medida posible, completas en la misma hoja del texto que le acompa a.



\begin{center}
    \section{Conclusiones}
\end{center}

        En esta secci n por lo general se describen los avances del trabajo, as  como una reflexi n con base en los resultados del trabajo realizado. Es com n tambi n mostrar ideas hacia trabajos futuros o aportaciones del trabajo realizado. Las conclusiones deben ser claras y concisas procurando enunciar los principales hallazgos o descubrimientos realizados.



\vspace{0.5cm}
% Para agregar o eliminar referencias, revisa el archivo "ref.bib", ah  encontrar s algunos ejemplos.
\nocite{*}



\begin{center}
    \section*{Autores}
\end{center}

    \textbf{Autor 1.} En  sta secci n se muestra una breve biograf a acad mica del autor, en no m s de 150 palabras. El nombre completo del autor de acuerdo a como se muestra al inicio del trabajo.

    \textbf{Autor 2.} En  sta secci n se muestra una breve biograf a acad mica del autor, en no m s de 150 palabras. El nombre completo del autor de acuerdo a como se muestra al inicio del trabajo.

    \textbf{Autor 3.} En  sta secci n se muestra una breve biograf a acad mica del autor, en no m s de 150 palabras. El nombre completo del autor de acuerdo a como se muestra al inicio del trabajo.
\end{document}